\documentclass[11pt,a4paper]{article}

\usepackage[utf8]{inputenc}
\usepackage[T1]{fontenc}
\usepackage{lmodern}
\usepackage[margin=1in]{geometry}

\usepackage{graphicx}
\usepackage{booktabs}
\usepackage{longtable}
\usepackage{float}
\usepackage{xcolor}
\usepackage{colortbl}
\usepackage{hyperref}
\usepackage{enumitem}
\usepackage{amsmath}

\definecolor{pass}{RGB}{46,204,113}
\definecolor{fail}{RGB}{231,76,60}
\definecolor{neutral}{RGB}{155,89,182}

\hypersetup{
    colorlinks=true,
    linkcolor=blue!60!black,
    urlcolor=blue!60!black,
    citecolor=blue!60!black,
}

\title{Code Coverage Agent: Ablation Study\\
\large Knowledge injection + structured execution for AI-generated tests}
\author{Mark Pollack}
\date{March 2026}

\begin{document}
\maketitle

\begin{abstract}
We evaluate the effect of domain knowledge injection and structured execution
on AI-generated test quality for Spring Boot projects. An AI agent (Claude Sonnet 4.6)
is given a Spring Boot project with production code and no tests, and asked to write
a comprehensive JUnit test suite. We compare five variants that progressively add
capability: prompt quality, domain knowledge files, and a two-phase explore-then-act
execution pattern. Each variant is evaluated by a four-tier automated judge cascade
measuring build success, coverage improvement, practice adherence, and similarity
to expert-written reference tests. On simple projects, a hardened prompt is the
dominant lever. Knowledge injection and structured execution show no measurable
benefit on trivial targets but are hypothesized to matter on complex, multi-layer
applications.
\end{abstract}

\tableofcontents
\newpage

%% ==================================================================
\section{Executive Summary}

Five variants progressively add capability: a naive prompt (control), a hardened
prompt with testing guidance (variant-a), curated domain knowledge (variant-b),
full knowledge base (variant-c), and a two-phase explore-then-act execution
pattern (variant-d).

Key findings from Sweep 001 (5 Spring Getting Started guides, N=1 per cell):

\begin{enumerate}
  \item A \textbf{hardened prompt} is the single biggest lever --- control fails 1/5 items; all other variants pass 100\%.
  \item All improved variants perform \textbf{similarly on simple projects} --- T3 adherence ranges 0.70--0.76, indistinguishable from noise at N=1.
  \item Two-phase execution costs \textbf{2$\times$} with no measurable benefit on simple targets.
  \item Coverage is a \textbf{ceiling}, not a differentiator --- 85--95\% across all variants.
  \item Golden similarity shows a \textbf{possible KB signal} on JPA and security items.
\end{enumerate}

\begin{table}[t]
\centering
\caption{Variant comparison (N=1 per cell, sorted by variant progression)}
\label{tab:variant-comparison}
\small
\begin{tabular}{lrlrlrll}
\toprule
 Agent:Model           &   N & Pass\%   &   Cov\% & T3   &   Eff & Total   & Per Item   \\
\midrule
 CC:Sonnet (naive)     &   5 & 100\%    &    95.6 & 0.62 &  0.88 & \$4.57  & \$0.91     \\
 CC:Sonnet (hardened)  &   5 & 100\%    &    85.7 & 0.80 &  0.94 & \$4.17  & \$0.83     \\
 CC:Sonnet (+3 KB)     &   5 & 100\%    &    95.6 & 0.70 &  0.84 & \$4.98  & \$1.00     \\
 CC:Sonnet (+full KB)  &   5 & 100\%    &    84.9 & 0.75 &  0.85 & \$5.13  & \$1.03     \\
 CC:Sonnet (two-phase) &   6 & 100\%    &    94.3 & 0.72 &  0.65 & \$16.48 & \$2.75     \\
 CC:Sonnet (forge)     &   1 & 100\%    &    95.3 & 0.95 &  0.17 & \$8.55  & \$8.55     \\
 CC:Sonnet (+SAE)      &   1 & 100\%    &    71.4 & 0.93 &  0.86 & \$0.92  & \$0.92     \\
 CC:Haiku              &   6 & 17\%     &    78.5 & 0.44 &  0.77 & \$2.88  & \$0.48     \\
 CC:Haiku (+SAE)       &   1 & 0\%      &    71.4 & 0.35 &  0.41 & \$0.48  & \$0.48     \\
 Loopy:Sonnet          &   1 & 100\%    &    94.6 & 0.73 &  0.86 & \$2.72  & \$2.72     \\
 Loopy:Haiku           &   6 & 33\%     &    71.4 & 0.49 &  0.86 & \$15.30 & \$2.55     \\
 Loopy:Qwen3           &   1 & 0\%      &     0   & ---  &  0.73 & \$5.01  & \$5.01     \\
\bottomrule
\end{tabular}
\end{table}


%% ==================================================================
\section{Experimental Design}

\subsection{Task}

Given a Spring Boot project with production code and no tests, generate a
comprehensive JUnit test suite targeting $\geq$80\% line coverage. The agent
operates in a fresh workspace with the project source, a prompt, and optionally
a set of domain knowledge files. It has full tool access: file reading/writing,
shell execution (compilation, test runs), and iterative repair.

The agent is Claude Sonnet 4.6. Single-phase variants (control through variant-c)
are invoked via \texttt{AgentClient} from the Spring AI Claude Agent library,
which provides a one-shot request--response pattern. Two-phase variant-d uses
\texttt{ClaudeSyncClient} from the Claude Agent SDK Java directly, because
\texttt{AgentClient} does not support session continuity --- variant-d needs
the explore phase's context preserved in memory for the act phase, which requires
two turns within the same Claude Code session.

In all cases, each invocation is self-contained: the agent reads the project,
writes tests, compiles, runs, checks coverage, and iterates until satisfied
or the timeout expires.

\subsection{Variants}

The five variants form a progressive ablation --- each adds one lever on top
of the previous:

\begin{table}[H]
\centering
\caption{Variant definitions}
\label{tab:variants}
\small
\begin{tabular}{@{}llllp{5cm}@{}}
\toprule
\textbf{Variant} & \textbf{Prompt} & \textbf{Knowledge} & \textbf{Phase} & \textbf{What changes} \\
\midrule
control   & v0-naive    & none          & single    & Minimal instructions: ``write tests for this project'' \\
variant-a & v1-hardened & none          & single    & Detailed testing guidance: slice annotations, assertion patterns, iteration protocol \\
variant-b & v2-with-kb  & 3 KB files    & single    & Adds 3 curated knowledge files (coverage fundamentals, JaCoCo patterns, Spring test slices) \\
variant-c & v2-with-kb  & full KB tree  & single    & Adds entire knowledge base (8 cheatsheets via index.md router) \\
variant-d & v3-explore + v3-act & full KB & two-phase & Two separate turns: first explore the project and KB, then write tests \\
\bottomrule
\end{tabular}
\end{table}

\subsection{Scoring: Four Dimensions}
\label{sec:scoring}

Each variant--item result is scored on four independent dimensions.
Scores are \textbf{not combined} into a single number --- they measure
different aspects of test quality.

\subsubsection{Line Coverage (JaCoCo)}

Standard JaCoCo line coverage percentage. Measured deterministically by running
\texttt{mvn clean test jacoco:report} on the agent's workspace after completion.
Reports both baseline (always 0\% since no tests exist initially) and final coverage.

\textbf{Interpretation}: Higher is better, but coverage alone does not measure
test quality --- a test that calls every method without assertions achieves high
coverage with no value.

\subsubsection{T3 Practice Adherence (LLM Judge)}

An LLM-based judge (Claude Sonnet 4.6) evaluates the generated test suite against
six criteria derived from Spring testing best practices:

\begin{enumerate}
  \item \textbf{Test slice selection} --- Does the test use the appropriate Spring test annotation (\texttt{@WebMvcTest}, \texttt{@DataJpaTest}, etc.) rather than \texttt{@SpringBootTest} everywhere?
  \item \textbf{Assertion quality} --- Are assertions specific and domain-aware (e.g., checking actual field values) rather than trivial (\texttt{assertNotNull})?
  \item \textbf{Error and edge case coverage} --- Does the suite test validation failures, missing entities, boundary conditions?
  \item \textbf{Domain-specific patterns} --- Are Spring-specific testing patterns used correctly (MockMvc for controllers, TestEntityManager for JPA)?
  \item \textbf{Coverage target selection} --- Does the suite focus on behaviorally significant code rather than trivial getters?
  \item \textbf{Version-aware patterns} --- Does the suite use APIs appropriate for the detected Spring Boot version (Boot 4.x vs 3.x)?
\end{enumerate}

Each criterion is scored 0.0--1.0 by the judge with supporting evidence.
\textbf{T3} is the mean across all six criteria. The rubric is identical for
all variants --- knowledge files inform the agent's test authorship, not the
judge's evaluation.

\textbf{Interpretation}: T3 measures \emph{how well} the tests follow expert
practices, independent of whether they compile or achieve coverage.

\subsubsection{Golden Test Similarity (AST Comparison)}

A deterministic comparison of the agent's test suite against expert-written
reference (``golden'') tests using JavaParser AST analysis. Five weighted
dimensions:

\begin{itemize}
  \item \textbf{Test method coverage} (weight 0.30) --- ratio of reference test methods that have a corresponding agent test
  \item \textbf{Annotation alignment} (weight 0.20) --- overlap in Spring test annotations used
  \item \textbf{Assertion style} (weight 0.20) --- ratio of assertion count (agent vs reference)
  \item \textbf{Import alignment} (weight 0.15) --- overlap in test imports (proxy for testing approach)
  \item \textbf{Injection pattern} (weight 0.15) --- overlap in injected dependencies (\texttt{@MockitoBean}, \texttt{@Autowired})
\end{itemize}

\textbf{Interpretation}: Golden similarity measures structural resemblance to
expert tests. A score of 1.0 means the agent produced tests structurally identical
to the reference. This is \emph{not} an LLM judge --- it is fully deterministic.

\subsubsection{Efficiency (Agent Behavior)}

Measures how efficiently the agent used its resources. Four weighted sub-metrics:

\begin{itemize}
  \item \textbf{Build errors} (weight 0.35) --- $1.0 - (\text{failed builds} / \text{total builds})$. An agent that compiles correctly on every attempt scores 1.0.
  \item \textbf{Cost efficiency} (weight 0.20) --- normalized inverse of total cost. Lower cost = higher score.
  \item \textbf{Recovery cycles} (weight 0.20) --- $1.0 - (\text{recovery iterations} / \text{total iterations})$. An agent that never needs to fix a broken build scores 1.0.
  \item \textbf{Tool utilization} (weight 0.25) --- ratio of productive tool calls to total tool calls.
\end{itemize}

\textbf{Interpretation}: Efficiency measures the agent's operational quality ---
did it work cleanly, or did it thrash through compile-fix-compile cycles?
Two agents can produce identical tests but differ dramatically in efficiency.

\subsection{Judge Cascade}

Judges are organized in a four-tier cascade. Lower tiers are binary gate checks;
higher tiers produce continuous scores. If a lower tier rejects (e.g., build fails),
higher tiers do not execute.

\begin{enumerate}[label=\textbf{Tier \arabic*:}, leftmargin=3cm]
  \item[Tier 0:] \textbf{BuildSuccessJudge} --- does the project compile and all tests pass? (binary gate)
  \item[Tier 1:] \textbf{CoveragePreservationJudge} --- no regression from baseline coverage. (binary gate)
  \item[Tier 2:] \textbf{CoverageImprovementJudge} + \textbf{GoldenTestComparisonJudge} --- how much did coverage improve? How similar are the tests to expert reference tests? (continuous scores)
  \item[Tier 3:] \textbf{TestQualityJudge} --- LLM-based practice adherence across 6 criteria. (continuous score)
\end{enumerate}

\subsection{Agent Execution and Capture}

Each agent invocation is captured via session log parsing (\texttt{SessionLogParser}
from the agent-journal library). Both \texttt{AgentClient} and \texttt{ClaudeSyncClient}
produce structured session logs that are parsed into a common \texttt{PhaseCapture}
record. The captured exhaust includes:

\begin{itemize}
  \item \textbf{Tool calls}: every file read, write, edit, glob, grep, and bash command
  \item \textbf{Token usage}: input tokens, output tokens, thinking tokens (extended thinking)
  \item \textbf{Turn count}: number of agent--environment round trips
  \item \textbf{Cost}: USD cost per invocation (API pricing)
  \item \textbf{Duration}: wall-clock time from prompt to final response
  \item \textbf{Phase count}: 1 for single-phase variants, 2 for variant-d (explore + act)
\end{itemize}

For two-phase variant-d, both phases run in the same \texttt{ClaudeSyncClient}
session --- \texttt{client.connect(explorePrompt)} followed by
\texttt{client.query(actPrompt)}. The explore phase's full context (tool
results, reasoning) is preserved in the session for the act phase.
Each phase's exhaust is captured separately via \texttt{SessionLogParser.parse()}.

\subsection{Dataset}

\textbf{Sweep 001}: 5 Spring Getting Started guides (bucket A --- simple).

\begin{table}[H]
\centering
\caption{Dataset items --- Sweep 001}
\label{tab:dataset}
\small
\begin{tabular}{@{}llrl@{}}
\toprule
\textbf{Item} & \textbf{Description} & \textbf{Classes} & \textbf{Complexity} \\
\midrule
gs-rest-service            & Single REST controller + record     & 2  & Trivial \\
gs-accessing-data-jpa      & JPA entity + repository             & 2  & Trivial \\
gs-securing-web            & MVC + Spring Security config        & 3  & Simple \\
gs-reactive-rest-service   & Reactive REST handler               & 2  & Trivial \\
gs-messaging-stomp-websocket & WebSocket STOMP controller        & 2  & Simple \\
\bottomrule
\end{tabular}
\end{table}

Each item has 1--3 golden reference test files written by domain experts,
used for the golden similarity comparison.

\textbf{Sweep 002} (planned): Spring Pet Clinic (bucket B --- multi-layer,
15+ production classes, form-based MVC, JPA entity relationships, validation,
Testcontainers integration tests, 17 golden reference test files).

%% ==================================================================
\section{Results: Getting Started Guides}

\subsection{Variant Comparison}

\begin{table}[t]
\centering
\caption{Judge cascade results per variant (tier 0--3 + golden)}
\label{tab:judge-cascade}
\small
\begin{tabular}{lllllllllllll}
\toprule
 Judge                 & CC:Sonnet (naive)   & CC:Sonnet (hardened)   & CC:Sonnet (+3 KB)   & CC:Sonnet (+full KB)   & CC:Sonnet (two-phase)   & CC:Sonnet (forge)   & CC:Sonnet (+SAE)   & CC:Haiku   & CC:Haiku (+SAE)   & Loopy:Sonnet   & Loopy:Haiku   & Loopy:Qwen3   \\
\midrule
 Build Success         & 100\%               & 100\%                  & 100\%               & 100\%                  & 100\%                   & 100\%               & 100\%              & 100\%      & 100\%             & 100\%          & 83\%          & 0\%           \\
 Coverage Preservation & 100\%               & 100\%                  & 100\%               & 100\%                  & 83\%                    & 0\%                 & 100\%              & 83\%       & 100\%             & 100\%          & 100\%         & nan\%         \\
 Coverage Improvement  & 100\%               & 100\%                  & 100\%               & 100\%                  & 83\%                    & 0\%                 & 100\%              & 83\%       & 100\%             & 100\%          & 100\%         & nan\%         \\
 Practice Adherence    & 0.62                & 0.80                   & 0.70                & 0.75                   & 0.72                    & 0.95                & 0.93               & 0.44       & 0.35              & 0.73           & 0.49          & nan           \\
 Golden Similarity     & nan                 & nan                    & nan                 & 0.68                   & 0.67                    & 0.58                & 0.36               & 0.49       & 0.43              & 0.88           & 0.60          & nan           \\
\bottomrule
\end{tabular}
\end{table}


All variants except control pass the binary gate checks (Tier 0--1) on all items.
Control fails on \texttt{gs-messaging-stomp-websocket} (build failure).
Coverage improvement is uniformly high (Tier 2) because the baseline is always
0\% and the agent consistently achieves 80--100\%.

The discriminating metrics are T3 practice adherence and golden similarity,
which show variance across variants.

\subsection{Per-Item Breakdown}

\begin{table}[t]
\centering
\caption{T3 practice adherence per item (higher is better)}
\label{tab:per-item-breakdown}
\small
\begin{tabular}{lllllrllrllrl}
\toprule
 Item                         & CC:Sonnet (naive)   & CC:Sonnet (hardened)   & CC:Sonnet (+3 KB)   & CC:Sonnet (+full KB)   &   CC:Sonnet (two-phase) & CC:Sonnet (forge)   & CC:Sonnet (+SAE)   &   CC:Haiku & CC:Haiku (+SAE)   & Loopy:Sonnet   &   Loopy:Haiku & Loopy:Qwen3   \\
\midrule
 gs-accessing-data-jpa        & 0.60                & 0.70                   & 0.65                & 0.68                   &                    0.73 & ---                 & ---                &       0.4  & ---               & 0.73           &          0.53 & ---           \\
 gs-messaging-stomp-websocket & 0.50                & 0.60                   & 0.50                & 0.50                   &                    0.5  & ---                 & ---                &       0.43 & ---               & ---            &          0.45 & ---           \\
 gs-reactive-rest-service     & 0.50                & 0.93                   & 0.70                & 0.90                   &                    0.73 & ---                 & ---                &       0.45 & ---               & ---            &          0.43 & ---           \\
 gs-rest-service              & 0.67                & 0.88                   & 0.85                & 0.88                   &                    0.88 & ---                 & 0.93               &       0.43 & 0.35              & ---            &          0.45 & nan           \\
 gs-securing-web              & 0.83                & 0.88                   & 0.78                & 0.78                   &                    0.67 & ---                 & ---                &       0.5  & ---               & ---            &          0.58 & ---           \\
 spring-petclinic             & ---                 & ---                    & ---                 & ---                    &                    0.8  & 0.95                & ---                &       0.45 & ---               & ---            &        nan    & ---           \\
\bottomrule
\end{tabular}
\end{table}


No single variant wins on all items. This is consistent with stochastic
variation at N=1 rather than a systematic advantage for any variant.

\begin{figure}[H]
\centering
\includegraphics[width=0.85\textwidth]{figures/per-item-bars.pdf}
\caption{T3 practice adherence by item and variant. Per-item variance is high;
no variant consistently dominates.}
\label{fig:per-item-bars}
\end{figure}

\subsection{Growth Story}

\begin{figure}[H]
\centering
\includegraphics[width=0.85\textwidth]{figures/growth-story.pdf}
\caption{Score progression as levers are added. The jump from control to
variant-a (hardened prompt) is the largest improvement. Subsequent levers
show diminishing returns on simple projects.}
\label{fig:growth-story}
\end{figure}

\subsection{Efficiency and Agent Behavior}

\begin{table}[t]
\centering
\caption{Efficiency breakdown per variant (higher is better, except tokens/duration)}
\label{tab:efficiency-breakdown}
\small
\begin{tabular}{lrrrrlll}
\toprule
 Agent:Model           &   Build Err &   Cost Eff &   Recovery &   Composite & Avg Out Tok   & Avg Think Tok   & Avg Duration   \\
\midrule
 CC:Sonnet (naive)     &        0.9  &       0.82 &       0.9  &        0.88 & 11,086        & 2,032           & 221s           \\
 CC:Sonnet (hardened)  &        0.97 &       0.83 &       0.97 &        0.94 & 9,869         & 1,558           & 216s           \\
 CC:Sonnet (+3 KB)     &        0.85 &       0.8  &       0.85 &        0.84 & 11,894        & 2,266           & 237s           \\
 CC:Sonnet (+full KB)  &        0.88 &       0.79 &       0.88 &        0.85 & 11,203        & 1,985           & 334s           \\
 CC:Sonnet (two-phase) &        0.65 &       0.55 &       0.75 &        0.65 & 26,436        & 4,866           & 508s           \\
 CC:Sonnet (forge)     &        0    &       0    &       0.62 &        0.17 & 60,528        & 6,854           & 1313s          \\
 CC:Sonnet (+SAE)      &        0.88 &       0.82 &       0.88 &        0.86 & 8,747         & 1,462           & 313s           \\
 CC:Haiku              &        0.67 &       0.9  &       0.81 &        0.77 & 25,472        & 3,752           & 294s           \\
 CC:Haiku (+SAE)       &        0    &       0.9  &       0.62 &        0.41 & 25,680        & 3,676           & 332s           \\
 Loopy:Sonnet          &        1    &       0.46 &       1    &        0.86 & 10,237        & 0               & 227s           \\
 Loopy:Haiku           &        1    &       0.49 &       1    &        0.86 & 26,448        & 0               & 330s           \\
 Loopy:Qwen3           &        1    &       0    &       1    &        0.73 & 24,486        & 0               & 1419s          \\
\bottomrule
\end{tabular}
\end{table}


Efficiency reveals a different story than quality scores. All single-phase
variants achieve similar quality, but their operational behavior differs:

\begin{itemize}
  \item \textbf{Build error rate}: variant-a scores highest (1.0 on most items) ---
  the hardened prompt's compilation-first iteration protocol reduces wasted cycles.
  Control and variant-d both show more build failures.

  \item \textbf{Cost efficiency}: Single-phase variants cost \$4--5 per sweep.
  Variant-d costs \$8.39 (2$\times$) because two Claude sessions are invoked
  per item. The explore phase adds context but also adds cost.

  \item \textbf{Token usage}: Variant-d generates the most output tokens
  (explore phase produces analysis text that is never directly used in tests).
  Knowledge-injected variants (b, c) use slightly more input tokens
  (reading KB files) but similar output tokens to variant-a.

  \item \textbf{Duration}: Single-phase variants complete in 2--5 minutes per item.
  Variant-d takes 5--10 minutes. Pet Clinic smoke test took 17 minutes for a single
  variant-a run --- complex projects will need longer timeouts.
\end{itemize}

\subsection{Cost Analysis}

\begin{table}[t]
\centering
\caption{Cost and token breakdown per variant}
\label{tab:cost-analysis}
\small
\begin{tabular}{lllllll}
\toprule
 Agent:Model           & Total   & Per Item   & In Tok     & Out Tok   & Think Tok   & Duration   \\
\midrule
 CC:Sonnet (naive)     & \$4.57  & \$0.91     & 109        & 55,430    & 10,160      & 1513s      \\
 CC:Sonnet (hardened)  & \$4.17  & \$0.83     & 137        & 49,343    & 7,791       & 1481s      \\
 CC:Sonnet (+3 KB)     & \$4.98  & \$1.00     & 137        & 59,472    & 11,328      & 1576s      \\
 CC:Sonnet (+full KB)  & \$5.13  & \$1.03     & 1,119      & 56,017    & 9,924       & 1671s      \\
 CC:Sonnet (two-phase) & \$16.48 & \$2.75     & 6,953      & 158,614   & 29,197      & 3046s      \\
 CC:Sonnet (forge)     & \$8.55  & \$8.55     & 96         & 60,528    & 6,854       & 1435s      \\
 CC:Sonnet (+SAE)      & \$0.92  & \$0.92     & 33         & 8,747     & 1,462       & 391s       \\
 CC:Haiku              & \$2.88  & \$0.48     & 1,905      & 152,835   & 22,509      & 2271s      \\
 CC:Haiku (+SAE)       & \$0.48  & \$0.48     & 317        & 25,680    & 3,676       & 370s       \\
 Loopy:Sonnet          & \$2.72  & \$2.72     & 854,353    & 10,237    & 0           & 314s       \\
 Loopy:Haiku           & \$15.30 & \$2.55     & 18,336,594 & 158,691   & 0           & 2470s      \\
 Loopy:Qwen3           & \$5.01  & \$5.01     & 1,547,477  & 24,486    & 0           & 1421s      \\
\bottomrule
\end{tabular}
\end{table}


\begin{figure}[H]
\centering
\includegraphics[width=0.7\textwidth]{figures/cost-quality.pdf}
\caption{Cost vs.\ quality trade-off. Variant-d (two-phase) costs 2$\times$ without
measurable quality improvement on simple targets. The efficient frontier is variant-a
(best quality-per-dollar).}
\label{fig:cost-quality}
\end{figure}

%% ==================================================================
\section{Results: Pet Clinic}

\textit{Sweep 002 in progress --- results pending.}

Smoke test (variant-a on spring-petclinic):
\begin{itemize}
  \item Line coverage: 93.6\% (277/296 lines)
  \item T3 practice adherence: 0.87
  \item Golden similarity: 0.66
  \item Cost: \$4.19
  \item Duration: 17 minutes (vs 2--3 minutes for guides)
  \item 60 test methods generated (reference has 59)
\end{itemize}

Full sweep (variant-a, variant-b, variant-d) will test whether knowledge injection
and structured execution provide measurable benefit on a complex, multi-layer project.

%% ==================================================================
\section{Analysis}

\subsection{Lever Ordering}

On simple projects, the dominant lever is prompt quality. The progression is:
\begin{enumerate}
  \item Hardened prompt ($+$0.10 T3 over naive baseline)
  \item Knowledge injection ($+$0.00--0.05 T3, within noise at N=1)
  \item Two-phase execution ($+$0.00 T3, $+$100\% cost)
\end{enumerate}

The hypothesis is that this ordering changes on complex projects where the
model cannot infer domain-specific testing patterns from context alone.

\subsection{Coverage as a Ceiling}

Coverage percentage is not useful for variant comparison on simple projects.
Three of five items show identical coverage across all variants (85--95\%).
The projects are too simple --- a single REST controller can be fully covered
by 2--3 test methods. Coverage becomes more discriminating on larger projects
(the Pet Clinic smoke test produced 93.6\% from significantly more code).

\subsection{Efficiency vs.\ Quality}

The efficiency dimension reveals that agents can produce equivalent test quality
through very different execution paths. Variant-a achieves the best
quality-per-dollar: highest efficiency score (0.94) with competitive T3 (0.71)
at the lowest cost (\$4.06). Variant-d spends 2$\times$ and gets no measurable
quality improvement on simple targets.

This suggests that on trivial tasks, additional ``thinking time'' (explore phase)
and ``reference material'' (knowledge base) are wasted effort --- the model already
knows enough. The question is where the complexity threshold lies.

\subsection{Limitations}

\begin{itemize}
  \item \textbf{N=1 per cell}: Every score is a single observation. Differences of $\sim$0.05 are indistinguishable from noise.
  \item \textbf{Simple dataset}: All items are single-layer, 1--3 class projects. Deliberately below the complexity threshold where knowledge should matter.
  \item \textbf{Same model}: All variants use Claude Sonnet 4.6. Cross-model comparison not yet tested.
  \item \textbf{No repeated trials}: Stochastic variation in LLM outputs means a single run may not represent typical behavior.
  \item \textbf{T3 is LLM-judged}: The practice adherence score comes from an LLM, which introduces its own noise. Golden similarity (AST-based) is deterministic but measures structure, not behavior.
\end{itemize}

%% ==================================================================
\section{Discussion}

The central thesis is:

\[
\textit{knowledge} + \textit{structured execution} > \textit{model alone}
\]

Sweep 001 on simple projects provides the \textbf{null condition}: when the task
is trivial, the model alone (with a good prompt) is sufficient. Knowledge injection
adds cost without measurable benefit. The two-phase pattern adds both cost and
complexity.

Sweep 001 is necessary but not sufficient to validate the thesis. The interesting
question is whether results change on complex targets. Pet Clinic (15+ classes,
4 architectural layers, form-based MVC, JPA entity relationships, validation,
Testcontainers) is designed to test exactly this.

If knowledge variants outperform on Pet Clinic but not on simple guides, the thesis
holds: \textbf{knowledge investment is justified when task complexity exceeds what
the model can infer from context alone}. If they don't, we learn that model
capability has grown past the point where curated knowledge matters --- at least
for Spring Boot testing.

%% ==================================================================
\section{Conclusion}

A hardened prompt is the minimum viable intervention for AI-generated test quality.
On simple Spring Boot projects, additional knowledge injection and structured
execution provide no measurable improvement at 2$\times$ the cost.

The experiment infrastructure --- automated judge cascade, golden test comparison,
efficiency tracking, and agent exhaust capture --- is validated and ready for
harder targets. Next: Spring Pet Clinic (complex target) and cross-model comparison
(knowledge on a cheaper model vs.\ a more capable model without knowledge).

\end{document}
